\chapter*{Prefácio}
\addcontentsline{toc}{chapter}{Prefácio}

Estas anotações acompanham a especialização em Engenharia e Administração de
Sistemas de Banco de Dados da Unicamp. A organização abaixo segue os quatro
blocos de formação e as ementas publicados na
\href{https://wordpress.ft.unicamp.br/espbd/?page_id=12}{página oficial de disciplinas do curso}.
Cada uma das 16 disciplinas tem carga horária de 24 horas.

\section*{Engenharia de Dados}

\begin{itemize}
	\item \textbf{CET 623 --- Artificial Intelligence Data Management:}
	uso de inteligência artificial para gerir o ciclo de vida dos dados e
	automatizar processos de validação, classificação e decisão.
	\item \textbf{CET 615 --- Fundamentos de Sistemas de Banco de Dados:}
	modelagem, estruturas de dados, restrições, índices, consultas SQL e
	programação no banco de dados.
	\item \textbf{CET 602 --- Engenharia de Requisitos de Banco de Dados:}
	levantamento, análise, especificação e validação de requisitos de dados,
	funcionais e não funcionais.
	\item \textbf{CET 603 --- Projeto de Bancos de Dados:}
	projeto conceitual, lógico e físico orientado por processos de negócio,
	qualidade de dados e metadados.
	\item \textbf{CET 616 --- Arquiteturas de Alto Desempenho para BD:}
	sistemas distribuídos, alta disponibilidade, infraestrutura e avaliação
	de desempenho de bancos de dados.
	\item \textbf{CET 619 --- BD Aplicados à Internet das Coisas (IoT):}
	armazenamento e plataformas de dados para objetos inteligentes e
	aplicações de Internet das Coisas.
	\item \textbf{CET 620 --- Bancos de Dados Não-relacionais:}
	modelos NoSQL de grafos, documentos, chave-valor e colunas, com práticas
	em Redis e MongoDB.
\end{itemize}

\section*{Administração de Banco de Dados (DBA)}

\begin{itemize}
	\item \textbf{CET 604 --- Programação e Otimização em SQL:}
	análise de planos de execução, melhoria de consultas e acompanhamento de
	métricas de desempenho.
	\item \textbf{CET 607 --- Segurança de Informações:}
	princípios e políticas de segurança, criptografia, gestão de riscos,
	continuidade e aspectos legais como a LGPD.
	\item \textbf{CET 617 --- Políticas de Administração de Bancos de Dados:}
	administração de instâncias, usuários, privilégios, armazenamento,
	memória, desempenho e alertas.
	\item \textbf{CET 624 --- Data Storage Management:}
	arquiteturas de armazenamento local e em nuvem, escalabilidade,
	disponibilidade e recuperação de desastres.
\end{itemize}

\section*{Ciência de Dados (Análise de Dados)}

\begin{itemize}
	\item \textbf{CET 622 --- Data Lake e Data Warehouse:}
	arquiteturas analíticas, modelagem dimensional, OLAP, ETL e lakehouse.
	\item \textbf{CET 611 --- Mineração de Dados:}
	preparação de dados, estatística e métodos de classificação, regressão,
	agrupamento e avaliação.
	\item \textbf{CET 612 --- Visualização de Informações:}
	representação de diferentes tipos de dados e aplicação de técnicas de
	visualização em análises e inteligência de negócios.
	\item \textbf{CET 621 --- Aprendizado de Máquina na Análise de Dados:}
	aprendizado supervisionado e não supervisionado, predição, avaliação de
	modelos e algoritmos como k-NN e redes neurais.
\end{itemize}

\section*{Arquitetura e Gestão de Dados}

\begin{itemize}
	\item \textbf{CET 618 --- Gestão e Governança de Dados:}
	alinhamento entre arquitetura, processos e governança, com atenção ao
	ciclo de vida, à qualidade e à maturidade da gestão de dados.
\end{itemize}

Ao final do curso, o Trabalho de Conclusão de Curso (TCC) aplica os
conhecimentos adquiridos em um módulo de 64 horas. Com as disciplinas, a
carga horária total informada pela Unicamp é de 448 horas.
